\InputIfFileExists{chapters/toc}{}{}


% how to read
\label{sec:key}\SectionBand{页面符号说明}{KEY TO THE NOTES}

\newcommand\legend[2]{\par\noindent\begin{minipage}[t]{46mm}\fontsize{10.2}{13}\selectfont\vspace{0pt}\raggedright#1\end{minipage}%
  \hfill\begin{minipage}[t]{72mm}\vspace{0pt}\fontsize{8.8}{13}\selectfont #2\end{minipage}\par\vspace{8pt}}

\legend{The rate \NWB{}{fell}{.}{fall}{/fɔːl/}{vi.}{下降}}
  {\textbf{新词}：实线。下面依次是音标、词性和本句中的意思；用的是变化形式时，先写原形。}
\legend{She \NS{}{runs}{}{新义}{run}{/rʌn/}{vt.}{经营} a shop.}
  {\textbf{新义}：学过的词出现新的意思，也用实线，下面标「新义」；词性也变了时标「新词性」。}
\legend{The rate \PH{went\ down}{go down}{下降}.}
  {\textbf{短语}：括线下写原形和意思。}
\legend{Its heart was \AGB{}{beating}{.}{beat}}
  {\textbf{本篇重复}：点线。形式有变化时，下面写原形。}
\legend{I felt a little \RV{}{nervous}{.}}
  {\textbf{复现词}：红色虚线。前面的文章学过，不再注释。}
\legend{the \PN{}{WHO}{}{/ˌdʌbljuː eɪtʃ ˈəʊ/}{n.}{世界卫生组织}\ \ a \GL{}{satellite}{}{/ˈsætəlaɪt/}{n.[C]}{卫星}}
  {\textbf{名称和超纲词}：灰色。不在词表内，了解即可。}
\legend{{\fontsize{8.8}{13}\selectfont n.[C]\enspace n.[U]\enspace n.[sing.]\\ n.[pl.]\enspace n.[C,U]\\ vt.\enspace vi.\enspace adj.\enspace adv.\\ prep.\enspace pron.\enspace conj.\\ det.\enspace num.\enspace int.}}
  {\textbf{词性}：按本句中的意思标注。n. 名词：[C] 可数，[U] 不可数，[sing.] 只用单数（常带 the），[pl.] 只用复数，[C,U] 可数和不可数都常用；专有名词（人名、地名、月份、星期等）只标 n.；vt. 及物动词，vi. 不及物动词，adj. 形容词，adv. 副词，prep. 介词，pron. 代词，conj. 连词，det. 限定词，num. 数词，int. 感叹词，modal v. 情态动词，abbr. 缩写。}
\legend{\SMK{y}She was half asleep.\EMK{y}{1}\ \ \SMK{g}He ran.\EMK{g}{1}}
  {\textbf{句子标记}：角括号括出的句子，在下方有同号的「译」（青色）或「语法」（橙色）注释。}
\legend{\leavevmode\llap{{\labelfont\bfseries\fontsize{8}{8}\selectfont\color{tealtext!70}3}\hspace{3mm}}Paragraph 3 starts here.}
  {\textbf{段落号}：每段开头的页边编号。}

\clearpage
